% !TeX program = xelatex
% UTF-8; compile with XeLaTeX from this directory or the repository root.
% Define \CarbohydrateDisordersInNotebook before inputting into a XePersian document.
% Source: Google Drive notebook/carbohydrate disorders, in filename order.
% pic01.jpg: 20261013_025829321.jpg (1Xsb3hfcaVYzIe8c9r3F-E2H9sGohzcgZ)
% pic02.jpg: 20261013_025836072.jpg (1pZ_fodSywdIc8r26sqrOw-H6iyc0fkHU)
% pic03.jpg: 20261013_025840130.jpg (1ZaUQQo0EY88S48MG5MRPaBJhR5ER8VQz)
\ifdefined\CarbohydrateDisordersInNotebook
  \let\CDfinish\relax
\else
  \def\CDfinish{\end{document}}
  \documentclass[12pt,a4paper]{article}
  \usepackage[margin=22mm]{geometry}
  \usepackage{amsmath,amssymb}
  \usepackage{graphicx}
  \usepackage{enumitem}
  \usepackage{xepersian}
  \IfFontExistsTF{Amiri}{\settextfont{Amiri}}{\settextfont{DejaVu Sans}}
  \setlatintextfont{DejaVu Serif}
  \setlength{\parindent}{0pt}
  \setlength{\parskip}{5pt}
  \setlength{\emergencystretch}{3em}
  \begin{document}
\fi
\providecommand{\CDmath}[1]{\lr{\(#1\)}}
\providecommand{\CDunclear}[1]{\textbf{[خوانش نامطمئن: #1]}}
\providecommand{\CDheader}[1]{%
  \section*{تصویر #1}
  \lr{Subject:} ۵۲ ـ اختلالات کربوهیدرات\par
  \lr{Year: \quad Month: \quad Date:} (خالی)\par
  شمارهٔ دست‌نویس صفحه: #1. شماره‌های چاپی کنار سطرها: \lr{1, 2, 3, 4, 5, 6, 7, 8, 9}.\par}

\begin{center}
{\Large\bfseries ۵۲ ـ اختلالات کربوهیدرات}
\end{center}
% Arrows are represented by the hierarchy and punctuation below.
% Uncertain readings are identified explicitly; originals are included at the end.

\CDheader{۱}
\begin{itemize}
\item تظاهرات اغلب بیماری‌های ذخیرهٔ گلیکوژن: \textbf{هایپوگلاسمی + هپاتومگالی}.
\item \textbf{انسولین}، تولید و تجمع گلیکوژن را تحریک می‌کند.
\item گلیکوژنولیز: تحریک توسط \textbf{اپی‌نفرین و گلوکاگون}.
\end{itemize}

\subsection*{انواع اختلالات ذخیرهٔ گلیکوژن}
\subsubsection*{تیپ \lr{0}}
گلیکوژن سنتاز دفیشنسی؛ کمبود \lr{Glycogen synthase}.
\begin{itemize}
\item ارگان درگیر: کبد.
\item سندرم بالینی: دورفیسم، هیپوگلاسمی ناشتا، کتوز.
\item نوزادان: تشنج، استفراغ هم هست.
\item \textbf{پروگنوز خوب}.
\end{itemize}

\subsubsection*{تیپ \lr{Ia}: \lr{Von-Gierke}}
کمبود آنزیم گلوکز ۶ ـ فسفاتاز.
\begin{itemize}
\item ارگان درگیر: کبد، کلیه، \lr{GI}، \textbf{پلاکت‌ها}.
\item هایپوگلاسمی، اسیدوز لاکتیک، \textbf{هپاتومگالی}، هایپوتونی، اسهال، خونریزی.
\item اگر سریع تشخیص داده نشود: \textbf{مرگ}.
\end{itemize}

\subsubsection*{تیپ \lr{Ib}}
کمبود گلوکز ۶ فسفات ترنسلوکاز.
\begin{itemize}
\item کبد، کلیه، \lr{GI}، \textbf{پلاکت، \lr{WBC}}.
\item مثل \lr{Ia} + نوتروپنی، \lr{IBD}، بیماری‌های \textbf{پریودنتال}.
\item \textbf{علائم \lr{IBD} شدید} دارد، \textbf{مرگ به علت عفونت در دههٔ دوم}.
\end{itemize}

\subsubsection*{تیپ \lr{II} (\lr{Pompe})}
لیزوزومال، \CDmath{\alpha} ـ گلوکوزیداز.
\begin{itemize}
\item همهٔ ارگان‌ها، عضلات مخطط، نورون‌ها.
\item ضعف شدید و پیشروندهٔ عضلات کمربند شانه‌ای، کاردیومگالی، هپاتومگالی، \lr{HF}، \lr{P-R} طولانی.
\item \lr{floppy baby}، ماکروگلوسی.
\item \textbf{پروگنوز بسیار بد}، \textbf{مرگ در سال اول}.
\end{itemize}

\subsubsection*{تیپ \lr{III} (\lr{Forbes})}
دبرنچینگ آنزیم؛ کبد، عضلات (\lr{IIIa})، یا کبد تنها (\lr{IIIb}).
\begin{itemize}
\item هایپوگلاسمی، کتون‌اوری، عقب‌ماندگی رشد، هپاتومگالی.
\item \textbf{پروگنوز بسیار بد}، مرگ در دههٔ اول به دلیل \textbf{نارسایی کبد}.
\end{itemize}

\clearpage
\CDheader{۲}
\subsubsection*{تیپ \lr{V} (\lr{McArdle})}
ماسل فسفریلاز، درگیری عضلات.
\begin{itemize}
\item خستگی، کرمپ عضلانی.
\item پروگنوز خوب، \textbf{طول عمر طبیعی}.
\end{itemize}

\subsubsection*{تیپ \lr{VI} (\lr{Hers})}
لیور فسفریلاز، درگیری کبد.
\begin{itemize}
\item هایپوگلاسمی خفیف، هپاتومگالی، کتون‌اوری.
\item \textbf{پروگنوز خوب}.
\end{itemize}

\subsubsection*{تیپ \lr{VII} (\lr{Tarui})}
ماسل فسفوفروکتوکیناز، درگیری عضلات.
\begin{itemize}
\item بالین شبیه تیپ \lr{V}.
\end{itemize}

\subsubsection*{تیپ \lr{VIII}}
فسفریلاز کیناز، درگیری کبد.
\begin{itemize}
\item شبیه تیپ \lr{III} بدون میوپاتی.
\item \textbf{پروگنوز خوب}.
\item وراثت وابسته به \lr{X}.
\end{itemize}

\begin{itemize}
\item بجز تیپ \lr{VIII}، بقیه وراثت \textbf{\lr{AR}} دارند.
\item \textbf{تیپ \lr{I} چربی بالا می‌رود}: هایپرتری‌گلیسریدمی.
\item تشخیص: اول تست \lr{DNA}، اگر نشد بیوپسی می‌کنیم.
\item تیپ \lr{II}: درمان با جایگزینی آنزیم.
\end{itemize}

\subsection*{گالاکتوزمی}
\begin{latin}
galactose-1-phosphate uridyltransferase
\end{latin}
\begin{itemize}
\item وراثت \textbf{\lr{AR}}.
\item علائم از \textbf{نوزادی} شروع می‌شود.
\item به دنبال شروع تغذیه با شیر مادر یا شیر گاو: \textbf{علائم نارسایی کبد}: هپاتومگالی، هایپربیلی‌روبینمی، اختلال انعقادی، هایپوگلاسمی.
\item علائم رنال توبولار: \textbf{اسیدوز}، گلوکوزاوری، آمینواسیداوری.
\item عدم تحمل غذا: اسهال، استفراغ.
\item علائم نورولوژیک: لتارژی، تشنج، تحریک‌پذیری.
\item \textbf{کاتاراکت}.
\end{itemize}

\clearpage
\CDheader{۳}
\subsection*{گالاکتوزمی (ادامه)}
\begin{itemize}
\item اگر درمان نشود: ظرف هفتهٔ اول می‌میرد.
\item در ریسک سپسیس با \textbf{\lr{E. coli}}.
\item علی‌رغم درمان: نارسایی زودرس تخمدان در خانم‌ها.
\item اختلالات یادگیری به درجات مختلف.
\item یافته‌های آزمایشگاهی: \textbf{هایپوگلاسمی، آلبومین‌اوری، گالاکتوزاوری}.
\item یادداشت حاشیه‌ای متصل به یافته‌های آزمایشگاهی: \CDunclear{«بیوپسی \lr{R/O} می‌کنه!»}.
\item تشخیص قطعی: نشان دادن فعالیت کم آنزیم.
\item \textbf{اختلال رنال توبولار}: \textbf{اسیدوز متابولیک هایپرکلرمیک نرمال آنیون گپ}.
\item درمان: حذف گالاکتوز.
\end{itemize}

\subsection*{گالاکتوکیناز دفیشنسی}
\begin{itemize}
\item \lr{AR}، تجمع گالاکتوز در مایعات بدن و تبدیلش به گالاکتیتول.
\item که به دلیل اثر اسموتیک: \textbf{کاتاراکت}، \CDmath{ICP\uparrow}.
\item \CDunclear{«اگر هردوژن»}: کاتاراکت نوزادی، قند خون نرمال.
\end{itemize}

\subsection*{عدم تحمل فروکتوز ارثی}
\begin{latin}
Hereditary fructose intolerance
\end{latin}
\begin{itemize}
\item مشابه گالاکتوزمی ولی با فروکتوز، سوکروز، سوربیتول.
\item تهوع، هایپوگلاسمی، نارسایی شدید کبد و کلیه.
\item حذف فروکتوز و سوکروز.
\end{itemize}

\subsection*{فروکتوزاوری}
\begin{itemize}
\item کمبود آنزیم فروکتوکیناز.
\item \textbf{عارضهٔ بالینی ندارد}.
\end{itemize}

\clearpage
\section*{تصاویر اصلی}
\IfFileExists{images/pic01.jpg}{\def\CDimagebase{images/}}{%
  \def\CDimagebase{pediatrics/ped-endocrinology/carbohydrate disorders/images/}}
\begin{center}
\includegraphics[width=\linewidth,height=0.88\textheight,keepaspectratio]{\CDimagebase pic01.jpg}
\end{center}
\clearpage
\begin{center}
\includegraphics[width=\linewidth,height=0.92\textheight,keepaspectratio]{\CDimagebase pic02.jpg}
\end{center}
\clearpage
\begin{center}
\includegraphics[width=\linewidth,height=0.92\textheight,keepaspectratio]{\CDimagebase pic03.jpg}
\end{center}
\CDfinish
